% ============================================================
%  TEMPLATE HANYA SOAL  (satu mata pelajaran per PDF)
%  Cover + Petunjuk Umum dalam satu halaman, lalu soal-soal.
%  Isi yang perlu diubah ditandai: % <-- UBAH
% ============================================================

\documentclass[11pt]{article}

% ===================== PACKAGES =====================
\usepackage[utf8]{inputenc}
\usepackage[T1]{fontenc}
\usepackage[a4paper,margin=2.5cm,top=2.8cm]{geometry}
\usepackage{amsmath,amssymb}
\usepackage{graphicx}
\usepackage{enumitem}
\usepackage{multicol}
\usepackage{xcolor}
\usepackage[most]{tcolorbox}
\usepackage{fancyhdr}
\usepackage{booktabs}
\usepackage{icomma}

% Tidak ada warna khusus -- hitam putih untuk cetak

% ===================== HEADER & FOOTER =====================
\pagestyle{fancy}\fancyhf{}
\lhead{\small NAMA UJIAN}                      % <-- UBAH
\rhead{\small Tahun}                           % <-- UBAH
\cfoot{\small\thepage}
\renewcommand{\headrulewidth}{0.4pt}

% ===================== KOTAK SOAL =====================
\newtcolorbox{soal}[1]{
  breakable,
  colback=white,
  colframe=black,
  colbacktitle=black,
  coltitle=white,
  fonttitle=\bfseries\sffamily,
  title=Nomor #1,
  boxrule=0.8pt, arc=0pt,
  left=3mm, right=3mm, top=2mm, bottom=2mm}

% ===================== PILIHAN GANDA =====================
\newenvironment{pil}{%
  \begin{enumerate}[label=(\Alph*),leftmargin=*,itemsep=1pt,topsep=3pt]}%
  {\end{enumerate}}

\newenvironment{pilB}{%
  \par\smallskip
  \begin{multicols}{2}
  \begin{enumerate}[label=(\Alph*),leftmargin=*,itemsep=1pt,topsep=3pt]}%
  {\end{enumerate}\end{multicols}}

% ===================== GAMBAR =====================
\newcommand{\gambar}[2][0.52\textwidth]{%
  \begin{center}\includegraphics[width=#1]{#2}\end{center}}

\linespread{1.05}

% ===================================================================
%                         ISI DOKUMEN
% ===================================================================
\begin{document}
\thispagestyle{empty}

% ===================== COVER + PETUNJUK (satu halaman) =====================

\begin{center}
  \fbox{\parbox{0.45\textwidth}{\centering\bfseries\large LEMBAR SOAL}}
\end{center}
\vspace{0.5cm}

\begin{center}
  {\LARGE\bfseries NAMA MATA PELAJARAN}        % <-- UBAH
\end{center}
\vspace{0.5cm}

\begin{center}
  \begin{tabular}{@{\hspace{2cm}}ll@{\hspace{0.4cm}:@{\hspace{0.4cm}}}l}
    Jenjang      & & [SMA / Madrasah Aliyah] \\[0.3em]  % <-- UBAH
    Hari/Tanggal & & [Hari, DD Bulan YYYY]   \\[0.3em]  % <-- UBAH
    Waktu        & & [X menit]               \\[0.3em]  % <-- UBAH
    Paket Soal   & & [A / B / --]            \\          % <-- UBAH / hapus baris ini jika tidak ada
  \end{tabular}
\end{center}

\vspace{0.6cm}
\noindent\rule{\linewidth}{0.6pt}
\vspace{0.3cm}

\noindent\textbf{\underline{PETUNJUK UMUM}}
\begin{enumerate}[leftmargin=*,itemsep=4pt,topsep=4pt]
  \item Anda \textbf{tidak} diperbolehkan menggunakan sumber-sumber eksternal,
        termasuk internet, \textit{large language model} seperti ChatGPT, Claude,
        LLaMA, Gemini, Meta AI, dan sebagainya. Anda sangat dianjurkan untuk
        mencoba mengerjakan sendiri terlebih dahulu karena evaluasi ini dibuat
        untuk \textbf{mengukur} tingkat kepahaman; nilai $<70$ mengindikasikan
        \textbf{tidak lulus}, dan jika sudah melewati \textit{deadline} maka
        dianggap tidak mengerjakan yang berarti \textbf{tidak lulus}.

  \item Terdapat 2 tahap pengerjaan, yaitu tahap \textbf{Try Out} dan tahap
        \textbf{Review}. Try Out bersifat \textit{open book}; Anda diperbolehkan
        membuka buku apapun selama bersifat \textit{offline} dengan rentang waktu
        2 jam pada tanggal \underline{\hspace{3cm}}. Setelah itu memasuki
        tahap Review dengan \textit{schedule} rentang tanggal
        \underline{\hspace{3cm}}--\underline{\hspace{3cm}}, di mana Anda
        diharapkan mengerjakan ulang dengan disertakan penjelasan/pembelaan
        mengapa revisi jawaban adalah pilihan tersebut, dan tetap mencantumkan
        penjelasan pada soal yang walaupun pada penilaian sudah benar.

  \item Anda dianjurkan untuk mencantumkan caranya walaupun tidak termasuk
        penilaian, mengingat sifatnya adalah sebagai bahan pembelajaran.
        Mohon bertanggung jawab atas pekerjaan Anda sendiri.

  \item Bacalah setiap soal dengan teliti dan pilih jawaban yang paling tepat.
        Terdapat soal yang jawabannya $\geq 1$ opsi.

  \item Soal berjumlah \underline{\;\;X\;\;} butir soal dengan bobot asli per
        soal adalah \underline{\;\;X\;\;} untuk mendapatkan nilai sempurna
        sebesar 100. Sebagian pada spesifikasi yang berlabel \textbf{[BONUS]};
        apabila dikerjakan dan jawaban tersebut tepat maka akan mendapat poin
        tambahan, dan apabila dikosongkan/salah tidak mengurangi poin. Khusus
        spesifikasi ini \textbf{diwajibkan} mencantumkan proses pengerjaan dan
        dikerjakan \textbf{hanya jika} keseluruhan soal terselesaikan.

  \item Silakan bertanya apabila terdapat soal yang ambigu, rancu, atau kurang
        spesifik selama itu tidak menjurus kepada jawaban soal tersebut. Setiap
        pertanyaan harus ada jawabannya dan setiap jawaban harus dapat
        dipertanggungjawabkan. Oleh karena itu terdapat kemungkinan apabila soal
        dianggap bonus/skor ($+$) apabila jawaban tersebut tidak dapat
        dipertanggungjawabkan ataupun tidak tercantum dalam pilihan jawaban.

  \item Isilah detail informasi berikut.\\[0.3em]
        \textbf{Nama}\hspace{0.45cm};\enspace\underline{\hspace{8cm}}
\end{enumerate}

\vspace{0.7cm}
\noindent\textit{Dengan menandatangani kolom di bawah ini, saya menyatakan telah membaca,
memahami, dan menyetujui seluruh peraturan yang tercantum di atas.}

\vspace{0.6cm}
\noindent\fbox{\parbox{5cm}{\vspace{2.2cm}\centering\small Tanda Tangan Peserta\vspace{0.3cm}}}

\clearpage

% ===================================================================
%  SOAL-SOAL
%  Hapus contoh di bawah dan isi dengan soal yang sebenarnya.
% ===================================================================

\begin{soal}{1}
[Teks soal nomor 1]
\begin{pilB}
  \item Opsi A \item Opsi B \item Opsi C \item Opsi D \item Opsi E
\end{pilB}
\end{soal}

\begin{soal}{2}
[Teks soal nomor 2 --- contoh dengan gambar]
\gambar[0.48\textwidth]{images/X/nama-gambar.png}
\begin{pil}
  \item Opsi A \item Opsi B \item Opsi C \item Opsi D \item Opsi E
\end{pil}
\end{soal}

\begin{soal}{3}
[Teks soal nomor 3 --- isian singkat, tanpa pilihan]
Nilai dari $\displaystyle\lim_{x\to 0}\frac{\sin 3x}{x}$ adalah \ldots
\end{soal}

% ... tambahkan soal-soal selanjutnya ...

% ===================================================================
%  DAFTAR JAWABAN SINGKAT
% ===================================================================
\clearpage
\begin{center}
  {\bfseries Daftar Jawaban Singkat}\\[0.6em]
  \begin{tabular}{c|*{10}{c}}
    \toprule
    No. & 1 & 2 & 3 & 4 & 5 & 6 & 7 & 8 & 9 & 10 \\
    \midrule
    Jwb & ? & ? & ? & ? & ? & ? & ? & ? & ? & ?  \\
    \bottomrule
  \end{tabular}
\end{center}

\end{document}
